\section{おわりに}
本稿では，モデルの再学習を行わずに，学習時に無い物体を検出し，確信度に応じた粒度で命名する粒度可変アノテーションを提案した．
提案手法は，YOLO11xのボックスにOWLv2の領域を加え，領域の特徴から求めた7806語の確率をWordNetの上位語へ集約して上位語から貪欲に下位語へ進み，各アノテーションと経路上の各単語に正しさの確率を付ける．
LVISの人手の注釈を学習に用いていないOWLv2による評価では，COCO外の物体の検出は56.0\%，命名は21.1\%であった．
同じ領域に対する命名の正しさは，確率が最大の1語を出力する場合の25.9\%から43.0\%に向上し，提案したスコアは期待校正誤差を0.145から0.028に下げ，全ての誤りを数えるAPを0.171から0.197に高めた．
これらの差は，ブートストラップによる信頼区間で確認した．
一方，改善の多くは粗い語による正解であり，LVISのカテゴリまで正しい出力のみを数える評価や，評価の語彙を事前に与えた場合には，集約の効果は無かった．
また，誤りの大部分は物体の一部や背景の領域に対するものであり，\add{ボックスの位置関係とWordNetの部分と全体の関係を用いてこれらを棄却する方法では，APの向上は0.003以下に止まった．}
今後は，領域が物体の全体であるかの判定，領域の特徴の分類能力の向上，およびGTと対応しない出力の人手による確認に取り組む．
